\documentclass[a4paper, 12pt]{article}

\usepackage{custom}


%--------------------VARIABILI--------------------
\def\lastversion{v0.1}
\def\title{Verbale interno}
\def\date{6 Maggio 2025}
%------------------------------------------------

\begin{document}

\primapagina

\begin{registromodifiche}
        \lastversion & 6 Maggio 2025 & Enrico Bianchi &  & Stesura verbale interno \\
        \hline 
\end{registromodifiche}

\tableofcontents

\newpage

\section{Registro presenze}
\begin{itemize}
    \item[] \textbf{Data}: 6 Maggio 2025
    \item[] \textbf{Ora inizio}:  10:00
    \item[] \textbf{Ora fine}: 11:00
    \item[] \textbf{Piattaforma}: Discord	
\end{itemize}

\begin{table}[H]
\centering
{\renewcommand{\arraystretch}{2}
\begin{tabularx}{\textwidth}{| X | X |}
    \hline
        \textbf{\large Componente} & 
        \textbf{\large Presenza} \\ 
    \hline 
    \hline
        Eghosa Matteo Igbinedion Osamwonyi&
        Presente \\
    \hline 
        Guirong Lan&
        Presente \\
    \hline 
        Enrico Bianchi&
        Presente \\
    \hline 
        Francesco Savio&
        Presente \\
    \hline 
        Marko Peric&
        Presente \\
    \hline 
        Pedro Leoni&
        Non presente \\
    \hline 

\end{tabularx}}
\end{table}

\newpage

\section{Verbale}

\subsection{Ordine del giorno}
\begin{itemize}
    \item Discussione sulle \glossario{issue} precedenti;
    \item Pianificazione delle attività da svolgere durante lo \glossario{Sprint} 14;
\end{itemize}    

\subsection{Decisioni prese}
    Il gruppo ha nominato Enrico Bianchi come responsabile per lo \glossario{sprint} 14.
    Le decisioni prese per questo sprint sono:
    \begin{itemize}
        \item Effettuare la stesura dei test di unità del backend.
        \item Effettuare la stesura dei test di integrazione del backend.
        \item Iniziare con la codifica di alcune componenti del frontend.
        \item Iniziare con la codifica della parte backend legata alle componenti frontend che si è deciso di implementare.
    \end{itemize}
    Per queste decisioni sono state pianificate le seguenti ore di lavoro:
    \begin{itemize}  
        \item 2 ore al responsabile
        \item 30 ore al programmatore.
        \item 16 ore al verificatore.
    \end{itemize}

\section{To Do/In progress}
\begin{tabularx}{\textwidth}{|c|C|C|C|}
    \hline
        \textbf{Issue} & \textbf{Descrizione} & \textbf{Incaricato} & \textbf{Scadenza}\\
    \hline 
        \#423 & Stesura del verbale interno & Enrico Bianchi & 7 Maggio 2025\\
    \hline
        \#424 & Aggiornamento del Piano di Progetto & Enrico Bianchi & 7 Maggio 2025\\
    \hline
        \#425 & Aggiornamento del Piano di Qualifica & Enrico Bianchi & 7 Maggio 2025\\
    \hline
        \#426 & Test di unità back-end & Pedro Leoni, Francesco Savio & 9 Maggio 2025\\
    \hline
        \#427 & Test di integrazione back-end & Pedro Leoni, Francesco Savio & 9 Maggio 2025\\
    \hline
        \#428 & Implementazione componente StandardPage e sottocomponenti associati & Marko Peric & 9 Maggio 2025\\
    \hline
        \#429 & Implementazione componente DatasetListPage e sottocomponenti associati & Guirong Lan & 9 Maggio 2025\\
    \hline
        \#430 & Implementazione back-end legata a DatasetListPage & Eghosa Matteo Igbinedion Osamwonyi & 9 Maggio 2025\\
    \hline
    \end{tabularx}

\end{document}